\chapter{Faithfully flat descent for modules and complexes}
\label{ch:flat-descent}

Let \(A\to B\) be a faithfully flat map of commutative rings. A
quasi-isomorphism can be checked after this base change because flatness
commutes with cohomology and faithfulness detects the zero module. This
chapter will also reconstruct an object from a strict cocycle and recover
an arbitrary unbounded complex from its augmented Amitsur diagram.

\section{The augmented complex and effective module descent}

\begin{definition}[The augmented Amitsur complex]\label{res:F01}
For an \(A\)-module \(M\), put \(C^{-1}(M)=M\) and
\(C^p(M)=B^{\otimes_A(p+1)}\otimes_A M\) for \(p\geq0\). The
map from degree \(-1\) inserts \(1\), and the map from degree \(p\)
is \(\delta=\sum_{i=0}^{p+1}(-1)^id^i\), where \(d^i\) inserts a
unit in position \(i\), counted from zero. This is a complex.
\end{definition}
\begin{proof}
Insertion in two ordered positions \(i<j\) can also be performed by
first inserting in position \(j-1\) and then in position \(i\).
The signs of the two resulting terms of the composite are opposite.
This cancels every term, including the two equal insertions after the
augmentation.
\end{proof}

\begin{theorem}[Exactness after faithfully flat descent]\label{res:F02}
The augmented Amitsur complex of every \(A\)-module is exact.
\end{theorem}
\begin{proof}
After tensoring on the left with \(B\), there is a contraction that
multiplies the new left factor with the first old factor:
\[
 h(b\otimes b_0\otimes\cdots\otimes b_p\otimes m)
   =bb_0\otimes b_1\otimes\cdots\otimes b_p\otimes m .
\]
In degree zero the right hand side is \(bb_0\otimes m\) in the
base-changed augmented term. We have \(hd^0=1\) and
\(hd^i=d^{i-1}h\) for \(i>0\). Therefore
\(h\delta+\delta h=1\), including \(h\delta=1\) on the augmented
term. By faithful flatness, homology of the original complex
must be zero.
\end{proof}

\begin{definition}[Strict module descent data]\label{res:F03}
Write \(B_2=B\otimes_A B\) and \(B_3=B^{\otimes_A3}\). A descent
datum on a \(B\)-module \(N\) is a \(B_2\)-linear isomorphism
\[
 \theta:N\otimes_A B\longrightarrow B\otimes_A N
\]
such that the pullback along multiplication \(B_2\to B\) is the
identity, and \(\theta_{23}\theta_{12}=\theta_{13}\) over \(B_3\).
Subscripts indicate the positions of the old and new \(N\)-factor.
A morphism is a \(B\)-linear map commuting with this isomorphism.
\end{definition}

\begin{theorem}[Effective descent for modules]\label{res:F04}
The functor \(M\mapsto B\otimes_A M\), with its canonical tensor
position isomorphism, is an equivalence onto modules with the
data of Result~\ref{res:F03}. The inverse sends \(N\) to
\[
 M=\{n\in N:\rho(n)=1\otimes n\},\qquad
 \rho(n)=\theta(n\otimes1).
\]
\end{theorem}
\begin{proof}
The map \(\rho\) is \(B\)-linear for the first factor of \(B\otimes_A N\).
The diagonal condition gives \(\mu\rho=1_N\), where \(\mu\) is
multiplication of scalars. If \(\rho(n)=\sum_j b_j\otimes n_j\),
evaluate the cocycle on \(n\otimes1\otimes1\) to get
\[
 (1\otimes\rho)\rho(n)=\sum_j b_j\otimes1\otimes n_j .
\]
Flatness identifies \(B\otimes_A M\) with the submodule of
\(B\otimes_A N\) whose elements \(z\) satisfy
\((1\otimes\rho)(z)=\sum b\otimes1\otimes n\), with
\(z=\sum b\otimes n\). The identity above places \(\rho(N)\)
in this equalizer. Conversely multiplying the first two factors in
the equalizer relation gives \(\rho(\mu z)=z\), using the \(B\)-linearity
of \(\rho\). Thus the natural map \(B\otimes_A M\to N\) is an isomorphism
with inverse \(\rho\) interpreted in the equalizer.
It respects the cocycle because for \(m\in M\) one has
\(\theta(m\otimes1)=1\otimes m\), and \(B_2\)-linearity determines
the isomorphism on all elements. In the canonical case
\(N=B\otimes_A M_0\), the equalizer is \(M_0\) by
Result~\ref{res:F02} in degrees \(-1,0\).
Both constructions are functorial since a map commuting with
\(\theta\) commutes with \(\rho\) and with \(\mu\).
\end{proof}

\begin{corollary}[Descent of finite projective modules]\label{res:F05}
For an \(A\)-module \(M\), if \(B\otimes_A M\) is finite projective
over \(B\), then \(M\) is finite projective over \(A\).
\end{corollary}
\begin{proof}
Express finitely many generators after base change using finitely many
elements of \(M\). The quotient by the module they generate has zero
base change, hence is zero. Choose a finite free surjection to \(M\).
Its kernel becomes finitely generated after base change, since the
surjection then splits and the kernel is a summand of a finite free
module. Apply the same finite-generation argument to this kernel.
Thus \(M\) is finitely presented.
It is flat because for any injection of \(A\)-modules, the kernel of
its tensor with \(M\) becomes zero over \(B\).
For completeness, a finitely presented flat module is projective.
At every localization at a prime, lifts of a residue-field basis
give a free surjection by Nakayama. Its finite kernel has zero
residue since flatness makes tensoring the sequence injective on
the left, so the kernel vanishes by Nakayama. A finite presentation
clears the denominators of this basis and its inverse on a
distinguished neighbourhood of the prime. The chosen finite free
surjection to \(M\) therefore splits over a finite distinguished
cover \(D(f_i)\). By finite presentation, clear denominators in the
local splittings to get maps \(\sigma_i:M\to A^n\) whose composites
with \(A^n\to M\) are \(f_i^{r_i}\mathrm{id}_M\) for sufficiently
large \(r_i\). These powers generate the unit ideal. A linear
combination of the \(\sigma_i\) gives a global splitting.
\end{proof}

\section{The unbounded descent calculation}

\begin{proposition}[Strict descent data for complexes]\label{res:F06}
Chain complexes of \(A\)-modules are equivalent to complexes of
\(B\)-modules equipped with the strict data of Result~\ref{res:F03}
by chain maps. The equivalence also identifies their localizations
at quasi-isomorphisms.
\end{proposition}
\begin{proof}
Apply Result~\ref{res:F04} in each degree. All differentials descend
as morphisms of the descent data. Their composites are zero by the
faithfulness of the equivalence. Tensoring with \(B\) commutes with
cohomology and reflects the zero module, so a map is a
quasi-isomorphism exactly when its base change is one.
The two functors and their comparison isomorphisms consequently
descend to the stated localizations.
\end{proof}

For a cochain complex \(X\) and an integer \(p\geq0\), put
\(C^{p,q}(X)=B^{\otimes_A(p+1)}\otimes_A X^q\).
Horizontal and vertical differentials commute before taking total
signs. Use the differential \(\delta+(-1)^pd_X\) on column \(p\),
and the product over \(p\) in each total degree.
This uses an actual product of modules, not a tensor product with
an already infinite product.

\begin{theorem}[Unbounded Amitsur totalization]\label{res:F07}
For every cochain complex \(X\) the natural augmentation is a
quasi-isomorphism
\[
       X\longrightarrow \operatorname{Tot}^{\Pi} C^{\bullet,\bullet}(X).
\]
The target is also the homotopy limit in \(D(A)\) of totalizations
truncated to the first finitely many columns.
\end{theorem}
\begin{proof}
For bounded-below \(X\), add the augmented column \(p=-1\) with
vertical sign \(-1\). The augmented total complex has acyclic rows
by Result~\ref{res:F02}. The lower bound on the vertical degree
and the finite number of columns in each total degree justify
the row filtration argument of \cite[B01--B04]{ha-core}. The
augmented total complex is the cone of the augmentation, with
column \(p=-1\) equal to \(X[1]\) and carrying differential \(-d_X\).
Its acyclicity proves the first assertion for bounded-below \(X\).

For general \(X\) put \(X_n=\tau_{\geq-n}X\) for \(n\geq0\), using
the quotient by boundaries at degree \(-n\), zero below and the
original modules above. The canonical tower maps \(X_{n+1}\to X_n\)
are degreewise surjective, the inverse limit is \(X\), and these
maps eventually become the identity in each fixed degree.
The augmentation maps for \(X_n\) form a map of towers.
For each column and each internal degree, tensoring with the indicated
finite tensor power of \(B\) commutes with this inverse limit because
the tower in that fixed degree is eventually constant.
Products commute with inverse limits, so the limit of the product
totalizations of the \(X_n\) is exactly the displayed product
totalization. Their transitions are degreewise surjective because
they are products of surjections.

For a tower \(Y_n\) with degreewise surjective transitions \(u_n\),
there is a short exact sequence of complexes
\[
 0\to\varprojlim_n Y_n\to\prod_{n\geq0}Y_n
 \xrightarrow{1-\mathrm{shift}}\prod_{n\geq0}Y_n\to0,
 \qquad (1-\mathrm{shift})(y)_n=y_n-u_n(y_{n+1}).
\]
Surjectivity follows by solving from \(n=0\) onwards. Products of
module complexes preserve quasi-isomorphisms because products of
surjections are surjective and kernels are products of kernels.
Apply this short exact sequence to the map of towers
\(X_n\to\operatorname{Tot}^{\Pi} C(X_n)\). The two product maps are
quasi-isomorphisms, so the long exact sequence gives the claim on
their inverse limits. This argument does not require \(B\otimes_A-\)
to preserve infinite products.

Finally truncate at horizontal column \(p=n\), by quotienting the
subcomplex of later columns. The projections between truncated
totalizations are surjective termwise, and the ordinary inverse
limit is the product totalization. The same short exact sequence
identifies that limit with the shifted cone of \(1-\mathrm{shift}\),
the sequential homotopy-limit model of Result~\ref{res:Q17}.
\end{proof}

\begin{corollary}[Descent of perfect-coefficient mapping complexes]\label{res:F08}
Let \(P\) be a bounded complex of finite projective \(A\)-modules
and \(X\) any complex. With \(B_p=B^{\otimes_A(p+1)}\) there is an
isomorphism in \(D(A)\)
\[
 \mathrm{RHom}_A(P,X)\simeq
 \operatorname{Tot}^{\Pi}_{p\geq0}
 \mathrm{RHom}_{B_p}(B_p\otimes_A P,B_p\otimes_A X),
\]
where the right hand side uses the cofaces of Result~\ref{res:F01}
on the compatible ordinary Hom models.
\end{corollary}
\begin{proof}
Bounded complexes of projectives are K-projective by the finite
filtration by degree tails. Thus the ordinary Hom complexes
compute each term. A finite projective module is a direct summand
of a finite free one. This gives
\(B_p\otimes_A\operatorname{Hom}_A(P^j,X^q)
=\operatorname{Hom}_{B_p}(B_p\otimes P^j,B_p\otimes X^q)\)
canonically, with equality here denoting the base-change isomorphism.
Each Hom degree uses only finitely many \(j\), so this remains
compatible with the Hom differential, multiplication and unit
on the base rings. Apply Result~\ref{res:F07} to
\(\operatorname{Hom}_A^\bullet(P,X)\).
\end{proof}

\begin{example}[Localization covers and the need for faithfulness]\label{res:F09}
If \(f_1,\ldots,f_s\) generate the unit ideal of \(A\), then
\(B=\prod_i A_{f_i}\) is faithfully flat and its Amitsur diagram
is the full ordered affine-intersection diagram of the cover by
\(D(f_i)\). In contrast, the flat map \(\mathbb Z\to\mathbb Q\)
does not detect the nonzero complex \(\mathbb Z/\ell\) in degree
zero for a prime \(\ell\).
\end{example}
\begin{proof}
The product is finite, hence a direct sum of flat modules as an
\(A\)-module. If \(m\) vanishes at each localization then it is
killed by a positive power of each \(f_i\). Those powers generate
the unit ideal, so \(m=0\). In particular base change reflects
zero modules, proving faithful flatness. The tensor product
distributes over the finite ring factors, and the tensor product
of the corresponding localizations is localization at the
product \(f_{i_0}\cdots f_{i_p}\). In the component indexed by a
longer ordered tuple, insertion of \(1\) restricts from the
intersection obtained by omitting that entry to the full intersection.
Finally \(\ell\) acts invertibly over \(\mathbb Q\), so tensoring
\(\mathbb Z/\ell\) with \(\mathbb Q\) gives zero.
\end{proof}

These are strict and explicit forms of descent. The general
homotopy-coherent formulation replaces ordinary compatibility by
a descent datum valued in a homotopy limit of higher categories.
Derived stacks and hypercovers use that framework
\cite[\S1.2.12]{hag}. The product calculation above already records
the higher Amitsur degrees and their total signs for affine
complexes.
